\chapter{Skin Biopsy}

\section*{Overview}
A skin biopsy removes a small sample of skin tissue for microscopic examination to diagnose skin conditions, including skin cancers, inflammatory dermatoses, or infections.

\section*{Alternative Names}
Dermatological biopsy, Punch biopsy, Shave biopsy

\section*{Principle}
A small sample of skin is removed by shaving the surface, punching a core, or excising the lesion and is submitted for histopathologic examination. Microscopic analysis provides a definitive diagnosis of skin conditions and tumors.
\section*{History}
Developed in the 19th century. The invention of the circular punch biopsy instrument in the 20th century standardized the collection of full-thickness skin cores.

\section*{Why It Is Done}
Ordered to evaluate suspicious moles, diagnose skin cancers (basal cell, squamous cell, melanoma), or identify chronic inflammatory skin disorders (psoriasis, eczema).

\section*{Is This a Primary or Secondary Test or Procedure}
Primary diagnostic test for suspected skin cancer and atypical dermatoses.

\section*{How It Is Performed}
The area is cleaned, and local anesthetic is injected. Depending on the technique: shave biopsy scrapes a superficial sample; punch biopsy punches a circular core; excisional biopsy cuts out the entire lesion.

\section*{Preparation}
Clean the skin. Tell the doctor if you take blood thinners or have a history of keloid scarring.

\section*{Contraindications and Precautions}
None. Active local skin infection is a relative contraindication (select a different site if possible).

\section*{Risks}
Minor bleeding, bruising, local infection, or scar formation (keloids).

\section*{Aftercare and Recovery}
Keep the biopsy site clean and covered with petroleum jelly and a bandage. If sutures were placed, they will be removed in 5-14 days.

\section*{Results and Interpretation}
The sample is analyzed by a dermatopathologist. Results describe the histopathology (e.g., benign nevus, basal cell carcinoma, or specific dermatitis).

\section*{Expected Results}
Normal epidermis and dermis; no cellular atypia, malignancy, or specific inflammatory pathology.

\section*{Follow-up}
Suture removal; review of pathology; surgical excision if malignant margins are positive.

\section*{References}
\begin{enumerate}
  \item MedlinePlus. Skin Biopsy. U.S. National Library of Medicine; 2025.
  \item Current Medical Diagnosis and Treatment. McGraw-Hill; 2025.
\end{enumerate}

\clearpage
